\subsection{Data Format}
The X-World model is trained on a meticulously curated dataset comprising a large scale of high-fidelity real-world driving sequences. These sequences are characterized by their diversity, covering a wide array of external environments, heterogeneous ego-vehicle behaviors, and complex multi-agent interactions. Each data sample constitutes a 10-second temporal segment and integrates the following multimodal data streams:

\begin{itemize}
    \item \textbf{Multiview Video Streams}: Synchronized video feeds from seven surrounding cameras.
    \item \textbf{Dynamic Object Trajectories}: Sequences of dynamic agents (e.g., vehicles, pedestrians) identified using a high-precision dynamic perception model.
    \item \textbf{Static Scene Elements}: Annotations of static infrastructure (e.g., lanes, traffic signs) obtained from a high-precision static perception model.
    \item \textbf{Textual Scene Descriptions}: Natural language descriptions of the driving scenario generated by the Vision Language Models (VLM).
\end{itemize}

% front_fisheye是否改成front_wide？ 1. 与front_narrow对应， 2. 更见名知义 -- camera布置图需要同步修改
The video data is recorded at a rate of 12 frames per second (FPS). Each frame provides a comprehensive 360-degree surround-view through seven distinct, calibrated camera perspectives: \textit{front\_narrow}, \textit{front\_fisheye}, \textit{front\_left}, \textit{front\_right}, \textit{rear\_left}, \textit{rear\_right}, and \textit{rear}. The precise spatial configuration and field-of-view overlap of these cameras are designed to ensure full coverage around the vehicle, as illustrated in Figure~\ref{fig:combined}(a).

% \begin{figure}[htbp]
%     \centering
%     \includegraphics[width=0.5\linewidth]{asset/pics/data/camera.png}
%     \caption{Spatial arrangement and fields of view of the seven camera sensors.}
%     \label{fig:camera_layout}
% \end{figure}
\begin{figure}[htbp]
    \centering
    \begin{subfigure}[b]{0.48\linewidth}
        \centering
        \includegraphics[width=\linewidth]{asset/pics/data/camera_layoutv2.pdf}
        % \caption{Spatial arrangement and fields of view of the seven camera sensors.}
        \label{fig:camera_layout}
    \end{subfigure}
    \hfill  % 添加水平间距
    \begin{subfigure}[b]{0.40\linewidth}
        \centering
    \includegraphics[width=\linewidth]{asset/pics/data/pie.pdf}
        \label{fig:longti_action}
    \end{subfigure}
    \caption{(a) Spatial layout of the seven camera sensors. (b) Distribution of ego longitudinal actions.}
    \label{fig:combined}
\end{figure}


\subsection{Video Captioning}

To enable fine-grained control and semantic understanding in X-World, we construct a large-scale video captioning pipeline tailored for autonomous driving scenes. Unlike generic video captions, our annotations focus on driving-specific attributes that are essential for controllable scene generation and downstream evaluation. 

\noindent\textbf{Captioning Schema.} Following our quantitative evaluation protocol, each video clip is annotated across four major dimensions:

\begin{itemize}
    \item \textbf{Macro Environment:} weather (sunny, cloudy, rainy, etc.), time of day (dawn, daytime, dusk, night), lighting condition, and driving environment (region type + road type).
    \item \textbf{Road Conditions:} surface condition (flat/bumpy), slope (uphill/downhill), and road state (dry/wet/puddles).
    \item \textbf{Traffic Infrastructure:} presence of lane markings, guardrails, traffic signs, traffic lights, buildings, vegetation, and special elements (bridges, construction zones, toll stations).
    \item \textbf{Traffic Density:} five-level scale from ``empty'' to ``congested''.
\end{itemize}


\noindent\textbf{Automated Pipeline.} Given the dataset, we adopt an automated approach using VLM. For each 10-second clip, we sampled a sequence of synchronized images from all 7 cameras in each 10-second fragment. The multi-view image sequence is fed into the model with a structured prompt that encodes the captioning schema.

\noindent\textbf{Example.} Below is a typical caption generated by our pipeline:

\begin{quote}
``The video captures a sunny daytime scene on a flat urban highway with sufficient sunlight. The road is lined with tall buildings and green vegetation, marked by clear white lane lines and roadside guardrails. A pedestrian overpass spans across the road in the distance. Traffic lights and signs are visible. Traffic density is moderate.''
\end{quote}

This structured, rule-guided approach ensures that every video clip in our dataset is paired with accurate, consistent, and semantically rich textual descriptions, laying the foundation for X-World's controllable generation capabilities.



\subsection{Auto Tagging}
To understand the natural distribution of data, enable more refined control over data distribution, and facilitate rapid data selection for small-scale feature validation, we developed a comprehensive, well-organized, and fine-grained three-level label taxonomy. Based on our task requirements, we defined four major categories of labels:

\begin{itemize}
    \item \textbf{Environmental Labels}: These describe overall scene-level characteristics and include 11 subcategories: Weather, Lighting, Road Surface Status, Road Surface Type, Road Curvature, Road Gradient, Road Structure, Road Type, Traffic Condition, Lane Clarity, and Lane Quantity Status. Each subcategory contains several fine-grained labels, resulting in a total of 50 third-level labels within the environmental category.
    \item \textbf{Static Labels}: Comprising 24 third-level labels, grouped into Road Markings, Lane Lines, Road Boundaries, Traffic Signs, Signal Applies To This Lane, Traffic Lights, and Static Obstacles.
    \item \textbf{Dynamic Labels}: Focus on five types describing traffic participants.
    \item \textbf{Ego-Vehicle Behavior Labels}: Include 21 third-level labels, mainly divided into Longitudinal, Lateral, Object Interaction, Scene Interaction, and Unreasonable Behavior.
\end{itemize}

The construction of this label taxonomy primarily relies on four information sources: (i) A high-precision dynamic perception network, which mainly serves dynamic labels; (ii) A high-precision static perception network, which mainly serves static labels; (iii) A robust online pose estimation system combined with vehicle pose information derived from high-precision sensors, which mainly serves ego-vehicle behavior labels; and (iv) the general-purpose VLM, which mainly serves environmental labels.

\subsection{Data Distribution}
We invested substantial computational resources to annotate the entire training dataset, enabling a comprehensive analysis of the natural data distribution and informed adjustments to the training set based on both statistical insights and model performance. Leveraging these labels, we conducted extensive iterative experiments, extracting valuable guidance for model training. As an example, Figure \ref{fig:combined}(b) illustrates the distribution of ego-vehicle longitudinal behaviors: the vast majority are normal driving (74.8\%), followed by stationary (21.0\%), while the remaining categories constitute a long tail. This analysis directly informs data collection—for instance, if the model exhibits poor performance on hard acceleration, we prioritize acquiring more such samples to enhance overall efficacy.

% \begin{figure}[htbp]
%     \centering
%     \includegraphics[width=0.5\linewidth]{asset/pics/data/Lontitunal.png}
%     \caption{Data Distribution of Ego Longitudinal Actions.}
%     \label{fig:longti_action}
% \end{figure}